\documentclass[11pt,letterpaper]{article}
\usepackage[T1]{fontenc}
\usepackage{lmodern,microtype}
\usepackage[margin=.85in,headheight=15pt]{geometry}
\usepackage{xcolor,booktabs,longtable,array,tabularx,enumitem,amsmath,graphicx,listings,tikz,fancyhdr,titlesec,hyperref,xurl}
\usetikzlibrary{arrows.meta,positioning,fit,shapes.geometric}
\definecolor{navy}{HTML}{16314D}\definecolor{teal}{HTML}{087F83}\definecolor{pale}{HTML}{EDF4F8}\definecolor{muted}{HTML}{536579}
\hypersetup{colorlinks=true,linkcolor=navy,urlcolor=teal,pdftitle={Agent Engineering Curriculum Integration}}
\titleformat{\section}{\Large\bfseries\color{navy}}{\thesection}{.7em}{}
\titleformat{\subsection}{\large\bfseries\color{teal}}{\thesubsection}{.7em}{}
\pagestyle{fancy}\fancyhf{}\fancyhead[L]{\small Agent Harness Engineering}\fancyhead[R]{\small Technical study / v1.0}\fancyfoot[R]{\thepage}
\setlength{\parindent}{0pt}\setlength{\parskip}{5pt}\renewcommand{\arraystretch}{1.18}
\setlist{nosep,leftmargin=*}\newcolumntype{Y}{>{\raggedright\arraybackslash}X}
\newcommand{\src}[1]{\expandafter\nolinkurl\expandafter{\detokenize{#1}}}
\newcommand{\evidence}[1]{\par\smallskip{\small\raggedright\color{muted}\textbf{Source anchor:} #1\par}}
\tikzset{n/.style={draw=navy,rounded corners=2pt,fill=pale,align=center,minimum height=9mm,text width=31mm,font=\small},g/.style={n,draw=teal,fill=teal!7},edge/.style={-Latex,thick,draw=navy},boundary/.style={draw=teal,dashed,rounded corners,inner sep=5pt}}
\lstset{basicstyle=\small\ttfamily,breaklines=true,frame=single,backgroundcolor=\color{pale},columns=fullflexible,showstringspaces=false}
\begin{document}
\begin{titlepage}
\pagenumbering{roman}
\vspace*{1.25in}
{\color{teal}\large AGENT HARNESS COURSE STUDY / 05\par}\vspace{.3in}
{\color{navy}\Huge\bfseries Agent Engineering Curriculum Integration\par}\vspace{.18in}
{\Large Source-to-lab crosswalk and proposed adoption plan\par}\vspace{.5in}
\rule{\linewidth}{1pt}\par
\textbf{Source:} Richmond Alake's Harness Engineering for AI Agents course export\par
\textbf{Version:} 1.0 \quad 23 November 2026\par
\textbf{Status:} Technical study; not a production authorization\par
\vfill
The model proposes; the harness validates, governs, executes and records.
\end{titlepage}
\pagenumbering{arabic}
\tableofcontents\clearpage
\section*{Reading note}
This standalone volume belongs to a six-document study of snapshot \src{8a5edab282632443}. Observed source behavior is distinguished from proposed controls. No live application, external provider or sandbox was independently executed. Test files are contract evidence, not independently verified passing results.

\section{Curriculum integration and implementation roadmap}
The target repository README and roadmap identify 16 labs, with L01--L02 described as reference implementations and L03--L16 as not implemented. This is the documented status, not a fresh test run. Retain the curriculum's five independent architecture layers: model, framework, harness, MCP and application governance. Its explicit prohibition on default shell/network/write tools, real GitHub changes, alert dismissal and AWS deployment overrides course examples that perform outward actions. Use synthetic fixtures and mock effects only. Live GitHub integration, when separately permitted by curriculum policy, is read-only.
\evidence{\url{https://github.com/Notoriousjayy/agent-engineering-curriculum/blob/main/README.md}; \url{https://github.com/Notoriousjayy/agent-engineering-curriculum/blob/main/docs/roadmap.md}.}
\subsection{Crosswalk and progressive laboratory specifications}
Each lab below is an extension proposal for its existing \src{labs/Lxx} location, not a claim of completion. Produce code, tests, threat/security notes, a Mermaid or equivalent diagram and real evidence.json with source SHA, versions, sanitized run ID and independent reviewer. Begin with existing lab contracts and preserve their acceptance criteria. The new criteria supplement rather than replace them.
\begin{longtable}{>{\raggedright\arraybackslash}p{.10\linewidth}>{\raggedright\arraybackslash}p{.34\linewidth}>{\raggedright\arraybackslash}p{.47\linewidth}}\toprule
Lab & Course anchor / goal & Proposed exercise and acceptance evidence \\\midrule\endhead
L01 / M01 & Part 1 stages 01, 05; distinguish five layers & Classify six scenarios; strict proposal and final schemas; reject malformed fields and unknown tools. Diagram model proposal vs authority. \\
L02 / M01 & Part 1 stages 05--06; bounded runner & Implement scripted model loop with read-only tool; stop at step/time budget; unknown tool denies; each dispatch has a trace. \\
L03 / M02 & Part 1 \src{langgraph_live.py}; parity & Implement explicit loop and graph adapter with same fixtures; compare normalized calls, results, termination and denials. \\
L04 / M02 & Part 3 \src{agent.py}; durable interrupts & SQLite checkpoint, mock effect interruption; restart before/after decision; verify no duplicate logical effect; unknown outcome reconciles. \\
L05 / M03 & Part 1 context, memory, skills, cache & Extract ports, context budgets and registry; compare baseline/context cost; prove cross-principal cache misses and no authority from memory. \\
L06 / M04 & Part 2 workspace MCP server & Build local server and host allowlist; validate schema and output limits; unlisted tools unavailable and invocation denied. \\
L07 / M05 & Part 2 workspace gateway analogy & Map MCP to read-only GitHub adapter; mock scopes then permitted live reads; deny cross-repo and mutation, redact credentials. \\
L08 / M06 & Part 1 semantic retrieval, Total Recall & Retrieve from synthetic repository docs with provenance, ACL and staleness checks; fixed corpus precision/recall report; midterm demonstration. \\
L09 / M07 & Part 2 policy, Part 3 System One & Threat model; inject malicious fixtures into mail-like docs and tool responses; zero unauthorized effect attempts executed; report detection misses. \\
L10 / M07 & Part 3 approval and action audit & Mock approval contract bound to actor and digest; reject tampering, expiry, replay, wrong actor and stale policy. No actual write API. \\
L11 / M08 & Part 1 parallel research; Part 3 graph & Delegate read-only subqueries, per-worker quotas, merge cited evidence; no worker privilege escalation, parent cancellation propagates. \\
L12 / M09 & Part 1 fair evaluation; Part 3 scorecard & Frozen task set, identical model/tools/budget; measure outcome, policy violations and latency; missing evidence blocks readiness. \\
L13 / M10 & Part 3 ledger and trajectories & Correlate model, tool, policy, effect and checkpoint spans; compute cost with configured rates; verify redaction and audit joins. \\
L14 / M11 & Part 1 sandbox and deployment & Offline CI builds and tests; pin dependencies, scan container, verify no secrets/egress and resource limits; no deployment. \\
L15 / M12 & Part 3 evidence graph adapted to CodeQL & Read-only alert investigation; collect source and alert provenance, cite evidence, never dismiss or modify; bounded investigation. \\
L16 / M12 & Part 3 learning gate and readiness & Independent reviewer tests attack/restart/isolation cases; optionally evaluate learned read-only procedure in shadow mode; BLOCKED on missing required evidence. \\
\bottomrule\end{longtable}
\section{Adoption and change control}
Keep L01 and L02 contracts intact. Add the crosswalk as a reference under \src{docs/reference/} only through a separate curriculum review. The curriculum README identifies L03--L16 as specified but not implemented; do not mark them complete until actual tests and evidence exist.
\begin{enumerate}
\item Inspect existing lab README, security notes, diagram and evidence file before proposing any change.
\item Prioritize L03/L04 orchestration parity and durability; then extract L05 ports and proceed through L06--L08 for the midterm.
\item Introduce adversarial fixtures and mock approval only in L09/L10; complete delegation and frozen benchmark before operational telemetry.
\item Finish L15 with read-only CodeQL investigation and L16 with independent evidence review. No real writes, alert dismissal or AWS deployment.
\end{enumerate}
The course's commercial and personal workflows supply design analogies, not authorization to reproduce their external effects. Reuse patterns, not private datasets, copied source or vendor-specific assumptions. Each lab retains its original acceptance criteria; the crosswalk adds proposed checks only.

\section{Source register and limitations}
\begin{itemize}
\item Richmond Alake, \textit{Harness Engineering for AI Agents}, supplied repository export,\newline snapshot \src{8a5edab282632443}; root README, Part 1--3 READMEs, code paths cited throughout. Workshop: \url{https://www.oreilly.com/live-events/harness-engineering-for-ai-agents/0642572381264}.
\item Target curriculum, README and roadmap, inspected 23 November 2026: \url{https://github.com/Notoriousjayy/agent-engineering-curriculum}.
\item Target documentation library, root README and AI technical-documents directory, inspected 23 November 2026: \url{https://github.com/Notoriousjayy/latex-docs}.
\end{itemize}
All figures are original logical reconstructions or proposed design diagrams; they are not copies of source graphics. Vendor-neutral controls, ADRs, acceptance tests and laboratory extensions are recommendations rather than claims about source implementation. Live dependencies, source test outcomes, provider claims, licensing/publication clearance and deployment configuration require separate validation.
\end{document}
